\honest{The Virtue of Narrowness}{Eliezer Yudkowsky}{2007}

I open with a line from my own ``Twelve Virtues'': ``more can be said about a single apple than about all the apples in the world.''

In their own trades, people know the value of narrow words. A mechanic does not lump a carburetor and a radiator together as ``car parts''; a hunter-gatherer knows a lion from a panther; a janitor does not wash the floor with window cleaner. Outside their trades, people ``often'' stretch a word as wide as it will go, since a general theory sounds loftier than an answer to a small question: better to explain human thought in general than how people solve a Rubik's Cube.

Here is a parable. The curious pick out a few pebbles and call them diamonds. Someone else wants to call every pebble a diamond, because leaving some out seems ``narrow-minded.'' \nb{In the parable the chosen pebbles really are diamonds, so the widener is wrong by construction. In a real dispute over a word, that is the question to be settled, and the post gives no way to settle it.} Even poets must be precise. Jealous, unconsummated love is not the love of a couple married for decades, and a flower for jealous love should have a heady scent, a bright color and thorns. Even when you shade meanings, you must keep track of which ones.

Biologists may explain living things and leave out the ``evolution'' of stars and of technology. Some ``unfortunate souls'' use one word for all three, as if a shared word made them one thing. \nb{The post names none of them. ``Stellar evolution'' is the astronomers' own term, with no claim of natural selection, and the word's general sense of development is older than the biological one.}

New Age gurus say ``Everything is connected to everything else,'' as if it were deep wisdom. It says nothing. A graph with every edge carries no more information than a graph with none; the useful graphs leave some things unconnected. People trying to sound profound compare this topic to that one until their graph is fully connected and useless. Detailed study shows how things are not alike, and lets you start ``subtracting edges.'' Good categories and good hypotheses exclude something. \nb{For hypotheses this is Karl Popper's criterion, not credited here.}

It was all right for Newton to explain gravity and not money or the heartbeat. \nb{Newton's gravity joined the fall of bodies on Earth to the orbits of the planets and the tides. It is remembered for a connection, not for a subtraction.} Sneering at narrowness recalls the ancient Greeks, for whom looking at things was manual labor, fit for slaves, and I quote Plato on the man who stares at a ceiling and learns nothing. \nb{Plato's passage is about astronomy done by gazing instead of by mathematics. It does not mention labor or slaves.}

Words too broad give ``something that isn't true and doesn't even make good poetry.'' And don't get me started on people who call Wikipedia an ``Artificial Intelligence,'' the invention of LSD a ``Singularity,'' or corporations ``superintelligent.''
